\documentclass[10pt]{article}
\usepackage[a4paper,margin=0.72in]{geometry}
\usepackage{amsmath,amssymb,booktabs,array}
\usepackage[T1]{fontenc}
\usepackage{lmodern}
\usepackage{microtype}
\setlength{\parindent}{0pt}
\setlength{\parskip}{4pt}
\begin{document}
\small

\begin{center}
{\Large\bf A Coordinate-Disjoint Cross-Holdout for the $c=1$ Prime-Shell Band}\\[5pt]
Liang Wang\\
School of Mathematics and Statistics, Huazhong University of Science and
Technology (HUST), Wuhan, China\\
September 4, 2026
\end{center}

\begin{abstract}
TPC-377 found that a finite $c=1$ prime-shell band retained its high-$Q$
failure support across a nested count ladder.  This paper tests whether the
support transfers to a fresh, coordinate-disjoint origin family.  We freeze
an affine grid, select three indices before reading any response, and evaluate
the endpoint counts $N=1024,2048$ at $Q=512,2048,8192$, giving 18 rows.  The
complete band profile is $(0,3,3)$ at both counts: 12 of 18 rows exceed the
spectral cap and none exceed the Schur cap.  Selected full-mode absolute
Rayleigh retention ranges from $0.93759972206138864$ to
$0.98046528117382914$.  This is finite response-blind support-transfer
evidence only; it does not prove origin or scale uniformity, an asymptotic
operator bound, an arithmetic estimate, or a twin-prime theorem.
\end{abstract}

\section{Question and claim boundary}

The recent TPC line has isolated a recurring finite signature: after
full-window square-energy normalization, the $c=1$ near-block band exceeds a
working spectral cap at the two high-$Q$ anchors while the $Q=512$ row remains
below it.  TPC-377 tested this signature on nested prefixes of established
origins.  The present question is deliberately narrower and more hostile:
does the same cap-support pattern occur on new origins whose coordinates are
disjoint from the earlier windows?

All statements in this paper concern the declared finite panel.  In
particular, a coordinate-disjoint sample is not an origin-uniform theorem,
and a shared left endpoint between two finite counts is not a growing
operator construction.  The normalization is recomputed at each scale.

\section{Finite object and exact identities}

For $I=[a,a+N-1]\cap\mathbb Z$, $p\in(Q,2Q]$, and $u,t\in I$, define
\[
 K_p(u,t)=p\left(\frac pQ\right)^2
 \frac{66^2}{66^2+(u-t)^2}
 \left({\bf1}_{p\mid(u-t)}-\frac1{p-1}\right)
 {\bf1}_{u\ne t}{\bf1}_{p\nmid u}{\bf1}_{p\nmid t}.
\]
The all-plus matrix and its full-window square-energy geometry are
\[
 A(u,t)=\sum_{Q<p\le2Q}K_p(u,t),\qquad
 G(u)=\sum_{t\in I}\sum_{Q<p\le2Q}K_p(u,t)^2,
 \qquad T(u,t)=\frac{A(u,t)}{\sqrt{G(u)G(t)}}.
\]
With $b(u)=\lfloor(u-a)/256\rfloor$, the inherited $c=1$ band is
\[
 B_1(u,t)=T(u,t){\bf1}_{|b(u)-b(t)|\le1},
 \qquad R_1=T-B_1.
\]
The geometry is a finite sum of nonnegative rational squares.  The exact
13-point anchor at $(a,Q)=(1100001,4)$ checks positive geometry and symmetry
by rational arithmetic.  The mask gives $T=B_1+R_1$ entrywise; hence, if
$Tv=\lambda v$ and $\|v\|_2=1$,
\[
 v^\mathsf{T}B_1v+v^\mathsf{T}R_1v
 =v^\mathsf{T}Tv=\lambda.
\]
These are exact finite identities, not uniform estimates.

\section{Predeclared cross-holdout protocol}

Before reading any response, we fix the affine candidate grid
\[
 a_j=1100001+401j,\qquad 0\le j<41,
\]
and select indices $(0,20,40)$, giving
$(1100001,1108021,1116041)$.  The largest current endpoint interval is
$[a,a+2047]$.  Exact integer endpoint comparisons show that all six current
intervals are disjoint from the largest declared TPC-376 and TPC-377
intervals.  At each new origin, the $N=1024$ interval is a prefix of the
$N=2048$ interval.

We freeze contiguous blocks of length 256, the $c=1$ mask, exponent one,
height 66, beta 2, the all-plus law, and $Q=(512,2048,8192)$.  The spectral
and Schur caps are respectively $0.64$ and $0.83$.  The complete Cartesian
panel has $3\times2\times3=18$ rows, and no row or count is selected using a
response, signed metric, or geometry score.  The extremal full mode is the
largest-absolute-eigenvalue mode, with the minimum mode resolving an exact
tie.

\section{Results}

Table~\ref{tab:profile} reports the range over the three fresh origins and
the number of spectral-cap failures.  The two rows in each entry are not
independent random samples; they are the fixed finite protocol above.

\begin{table}[ht]
\centering
\caption{Band spectral ranges and failure profile.}
\label{tab:profile}
\scriptsize
\begin{tabular}{c c c c c}
\toprule
$N$ & blocks & $Q=512$ & $Q=2048$ & $Q=8192$\\
\midrule
1024 & 4 & .60988488--.60989272 & .65205068--.65205804 &
.65334113--.65334662\\
2048 & 8 & .50283355--.50284232 & .66562664--.66564016 &
.66694192--.66694513\\
\midrule
\multicolumn{2}{c}{failure count} & $0/3$ & $3/3$ & $3/3$\\
\bottomrule
\end{tabular}
\end{table}

Thus the count-by-$Q$ profile is
\[
 (0,3,3)\quad\text{at }N=1024,\qquad
 (0,3,3)\quad\text{at }N=2048.
\]
There are 12 spectral-cap violations and zero Schur-cap violations.  The
profile agrees with TPC-377 on this new coordinate-disjoint panel.  The
agreement is a support statement: the displayed spectral ranges change
between counts, so no magnitude-stability law is inferred.

For the selected full eigenmode, the absolute band-Rayleigh retention over
all 18 rows satisfies
\[
 0.93759972206138864\le
 \frac{|v^\mathsf{T}B_1v|}{|\lambda|}
 \le0.98046528117382914,
\]
and the largest corresponding absolute tail fraction is
$0.062400277938610291$.  Eigen residual, norm, symmetry, Frobenius, and
Schur-envelope checks are recorded row by row in the canonical certificate.

\section{Independent audit}

The producer locks the TPC-377 source and certificate and writes a canonical
JSON certificate.  A separate checker does not import the TPC-378 producer:
it sieves primes through 20000, accumulates the shell in reverse order,
reconstructs the full and masked matrices, and recomputes all 18
eigensystems.  It also recomputes the exact rational anchor.  A 24-mutation
stress suite changes protocol, row, census, and firewall fields and requires
every mutation to be rejected.

Normal and optimized Python runs are required to return zero with empty
standard error and byte-identical summary lines.  The repository Bridge-B
then locks the package source, certificate, paper, PDF, log, and route notes
before repeating these checks.  The official Session-named Route-A and
Route-B evaluator files are absent, so Bridge-B is local fail-closed evidence
and not an official evaluator verdict.

\section{Limitations and route decision}

The finite transfer does not prove origin uniformity, window-scale
uniformity, spectral-magnitude stability, cross-block causality, source
validity of the normalization, a growing masked-operator bound, or a
source-uniform arithmetic $L^2$ estimate.  It pays no fixed-power credit:
\texttt{ARITHMETIC\_ADVANCE = NO}, \texttt{FIXED\_POWER\_CREDIT = 0}, and
\texttt{FULL\_GATE\_B = OPEN}.  There is no Route-B reassembly and no
twin-prime conclusion.

The strongest positive result is a response-blind profile transfer to three
fresh coordinate-disjoint origins at two endpoint scales.  The strongest
obstruction is that this transfer remains a threshold census with
scale-specific normalization and moving magnitudes.  The reusable structure
is the affine-grid selector together with exact interval separation and the
common band/tail Rayleigh audit.  The next finite question is
\texttt{TEST\_C1\_CROSSHOLDOUT\_LAW\_CONTROL}.

\section{Conclusion}

TPC-378 closes the finite origin/scale cross-holdout proposed by TPC-377:
the parent $(0,3,3)$ support pattern survives on the declared fresh panel.
It strengthens the empirical map of the finite model while leaving every
arithmetic and growing-operator gate open.

\end{document}
